
\documentclass{article}
\usepackage{colortbl}
\usepackage{makecell}
\usepackage{multirow}
\usepackage{supertabular}

\begin{document}

\newcounter{utterance}

\centering \large Interaction Transcript for game `clean\_up', experiment `0\_easy\_5obj\_lt', episode 2 with Apertus{-}v1.5{-}8B.
\vspace{24pt}

{ \footnotesize  \setcounter{utterance}{1}
\setlength{\tabcolsep}{0pt}
\begin{supertabular}{c@{$\;$}|p{.15\linewidth}@{}p{.15\linewidth}p{.15\linewidth}p{.15\linewidth}p{.15\linewidth}p{.15\linewidth}}
    \# & $\;$A & \multicolumn{4}{c}{Game Master} & $\;\:$B\\
    \hline

    \theutterance \stepcounter{utterance}  
    & & \multicolumn{4}{p{0.6\linewidth}}{
        \cellcolor[rgb]{0.9,0.9,0.9}{
            \makecell[{{p{\linewidth}}}]{
                \texttt{\tiny{[P1$\langle$GM]}}
                \texttt{Aš esu jūsų žaidimo vedėjas, o jūs žaidžiate bendrą žaidimą, kurio žaidimo lenta yra tokia lentelė:} \\
\\ 
\texttt{```} \\
\texttt{ 1234567} \\
\texttt{╔═══════╗} \\
\texttt{║◌◌◌◌◌◌◌║ 1} \\
\texttt{║◌◌◌◌◌◌◌║ 2} \\
\texttt{║◌TW◌◌┌─╢ 3} \\
\texttt{╟──┐◌◌│H║ 4} \\
\texttt{║C◌├──┼─╢ 5} \\
\texttt{║◌◌│I◌│◌║ 6} \\
\texttt{║◌◌│◌◌│◌║ 7} \\
\texttt{╚══╧══╧═╝} \\
\\ 
\texttt{```} \\
\\ 
\texttt{* Viršutiniame krašte x koordinatės didėja į dešinę, o dešiniajame krašte y koordinatės didėja žemyn.} \\
\texttt{* Tinklelyje atsitiktine tvarka išdėstyti šie objektai: 'W', 'I', 'T', 'C', 'H'.} \\
\\ 
\\ 
\texttt{Kitas žaidėjas mato kitą žaidimo lentos variantą, kuriame objektai išdėstyti skirtingose atsitiktinėse vietose. Jūs negalite matyti kito žaidėjo lentos, o jis negali matyti jūsų lentos.} \\
\\ 
\\ 
\texttt{**Tikslas:**} \\
\\ 
\texttt{Abu žaidėjai turi perkelti objektus savo fone taip, kad vienodi objektai atsidurtų tose pačiose koordinatėse. Turite bendrauti su kitu žaidėju ir susitarti dėl bendros tikslinės konfigūracijos.} \\
\\ 
\\ 
\texttt{**Taisyklės:**} \\
\\ 
\texttt{* Kiekvieno ėjimo metu galite išsiųsti lygiai vieną iš šių dviejų komandų:} \\
\texttt{ 1. `SAKYTI: <ŽINUTĖ>`: norėdami išsiųsti žinutę (viską iki kitos eilutės pertraukos) kitam žaidėjui. Aš ją persiųsiu jūsų partneriui.} \\
\texttt{2. `JUDĖTI: <OBJEKTAS>, (<X>, <Y>)`: perkelti objektą į naują poziciją, kur `<X>` yra stulpelis, o `<Y>` yra eilutė. Aš jums pranešiu, ar jūsų perkėlimas buvo sėkmingas, ar ne.} \\
\texttt{* Jei nesilaikysite formato arba siųsite kelias komandas iš karto, turėsiu jus nubausti.} \\
\texttt{* Jei abu žaidėjai sukaups daugiau nei 12 baudų, abu pralaimi žaidimą.} \\
\texttt{* Labai svarbu, kad bendrautumėte su kitu žaidėju dėl savo tikslinės būsenos! Vienintelis būdas, kuriuo galite perduoti savo strategiją kitam žaidėjui, yra komanda `SAKYTI: <ŽINUTĖ>`!} \\
\\ 
\\ 
\texttt{**Objektų perkėlimas**} \\
\\ 
\texttt{* Objektus galima perkelti tik į langelius, esančius tinklelio ribose. Tikslinis langelis turi būti tuščias, t. y. jame turi būti tik simbolis '◌'.} \\
\texttt{* Jei bandysite perkelti objektą į vietą, kuri nėra tuščia, arba bandysite jį perkelti už tinklelio ribų, turėsiu jus nubausti. Gausite dar vieną bandymą.} \\
\texttt{* Prieš atlikdami perkėlimą, dukart patikrinkite, ar tikslinė vieta yra tuščia ir ar joje nėra jokios raidės, rėmelio ar linijos!} \\
\\ 
\\ 
\texttt{**Žaidimo pabaiga**} \\
\\ 
\texttt{Jei manote, kad pasiekėte tikslą {-} sulygiavote visus objektus, galite paprašyti kito žaidėjo baigti žaidimą siųsdami `SAKYTI: baigta?`. Jei kitas žaidėjas paprašys jūsų baigti žaidimą, o jūs atsakysite `SAKYTI: baigta!`, žaidimas baigsis.} \\
\\ 
\texttt{Abu žaidėjai laimi, jei žaidimas baigiasi per 20 raundus, kur vienas raundas apibrėžiamas kaip du žaidėjai, kiekvienas išsiuntęs po vieną galiojančią komandą.} \\
\\ 
\texttt{**Taškų skaičiavimas:**} \\
\\ 
\texttt{Kuo arčiau vienodų objektų abiejose žaidimo lentose, tuo daugiau taškų gausite. Baudos mažina jūsų taškus. Ar pavyks pagerinti rekordą?} \\
\\ 
\\ 
\texttt{Išsiųskite žinutę kitam žaidėjui, kad pradėtumėte žaidimą!} \\
            }
        }
    }
    & \\ \\

    \theutterance \stepcounter{utterance}  
    & \multicolumn{4}{p{0.6\linewidth}}{
        \cellcolor[rgb]{0.9,0.9,0.9}{
            \makecell[{{p{\linewidth}}}]{
                \texttt{\tiny{[P1$\rangle$GM]}}
                \texttt{SAKYTI: norėčiau susitarti dėl galutinės konfigūracijos. Prašau parodyti savo viziją, kur būtų išdėstyti objektai W, I, T, C, H.} \\
            }
        }
    }
    & & \\ \\

    \theutterance \stepcounter{utterance}  
    & & & \multicolumn{2}{p{0.3\linewidth}}{
        \cellcolor[rgb]{0.9,0.9,0.9}{
            \makecell[{{p{\linewidth}}}]{
                \texttt{\tiny{[GM$|$GM]}}
                \texttt{SAKYTI: norėčiau susitarti dėl galutinės konfigūracijos. Prašau parodyti savo viziją, kur būtų išdėstyti objektai W, I, T, C, H.} \\
            }
        }
    }
    & & \\ \\

    \theutterance \stepcounter{utterance}  
    & & \multicolumn{4}{p{0.6\linewidth}}{
        \cellcolor[rgb]{0.9,0.9,0.9}{
            \makecell[{{p{\linewidth}}}]{
                \texttt{\tiny{[GM$\rangle$P2]}}
                \texttt{Aš esu jūsų žaidimo vedėjas, o jūs žaidžiate bendrą žaidimą, kurio žaidimo lenta yra tokia lentelė:} \\
\\ 
\texttt{```} \\
\texttt{ 1234567} \\
\texttt{╔═══════╗} \\
\texttt{║◌◌◌◌TH◌║ 1} \\
\texttt{║◌◌◌◌◌◌◌║ 2} \\
\texttt{║◌◌◌◌◌┌─╢ 3} \\
\texttt{╟──┐◌◌│◌║ 4} \\
\texttt{║I◌├──┼─╢ 5} \\
\texttt{║◌W│◌◌│C║ 6} \\
\texttt{║◌◌│◌◌│◌║ 7} \\
\texttt{╚══╧══╧═╝} \\
\\ 
\texttt{```} \\
\\ 
\texttt{* Viršutiniame krašte x koordinatės didėja į dešinę, o dešiniajame krašte y koordinatės didėja žemyn.} \\
\texttt{* Tinklelyje atsitiktine tvarka išdėstyti šie objektai: 'W', 'I', 'T', 'C', 'H'.} \\
\\ 
\\ 
\texttt{Kitas žaidėjas mato kitą žaidimo lentos variantą, kuriame objektai išdėstyti skirtingose atsitiktinėse vietose. Jūs negalite matyti kito žaidėjo lentos, o jis negali matyti jūsų lentos.} \\
\\ 
\\ 
\texttt{**Tikslas:**} \\
\\ 
\texttt{Abu žaidėjai turi perkelti objektus savo fone taip, kad vienodi objektai atsidurtų tose pačiose koordinatėse. Turite bendrauti su kitu žaidėju ir susitarti dėl bendros tikslinės konfigūracijos.} \\
\\ 
\\ 
\texttt{**Taisyklės:**} \\
\\ 
\texttt{* Kiekvieno ėjimo metu galite išsiųsti lygiai vieną iš šių dviejų komandų:} \\
\texttt{ 1. `SAKYTI: <ŽINUTĖ>`: norėdami išsiųsti žinutę (viską iki kitos eilutės pertraukos) kitam žaidėjui. Aš ją persiųsiu jūsų partneriui.} \\
\texttt{2. `JUDĖTI: <OBJEKTAS>, (<X>, <Y>)`: perkelti objektą į naują poziciją, kur `<X>` yra stulpelis, o `<Y>` yra eilutė. Aš jums pranešiu, ar jūsų perkėlimas buvo sėkmingas, ar ne.} \\
\texttt{* Jei nesilaikysite formato arba siųsite kelias komandas iš karto, turėsiu jus nubausti.} \\
\texttt{* Jei abu žaidėjai sukaups daugiau nei 12 baudų, abu pralaimi žaidimą.} \\
\texttt{* Labai svarbu, kad bendrautumėte su kitu žaidėju dėl savo tikslinės būsenos! Vienintelis būdas, kuriuo galite perduoti savo strategiją kitam žaidėjui, yra komanda `SAKYTI: <ŽINUTĖ>`!} \\
\\ 
\\ 
\texttt{**Objektų perkėlimas**} \\
\\ 
\texttt{* Objektus galima perkelti tik į langelius, esančius tinklelio ribose. Tikslinis langelis turi būti tuščias, t. y. jame turi būti tik simbolis '◌'.} \\
\texttt{* Jei bandysite perkelti objektą į vietą, kuri nėra tuščia, arba bandysite jį perkelti už tinklelio ribų, turėsiu jus nubausti. Gausite dar vieną bandymą.} \\
\texttt{* Prieš atlikdami perkėlimą, dukart patikrinkite, ar tikslinė vieta yra tuščia ir ar joje nėra jokios raidės, rėmelio ar linijos!} \\
\\ 
\\ 
\texttt{**Žaidimo pabaiga**} \\
\\ 
\texttt{Jei manote, kad pasiekėte tikslą {-} sulygiavote visus objektus, galite paprašyti kito žaidėjo baigti žaidimą siųsdami `SAKYTI: baigta?`. Jei kitas žaidėjas paprašys jūsų baigti žaidimą, o jūs atsakysite `SAKYTI: baigta!`, žaidimas baigsis.} \\
\\ 
\texttt{Abu žaidėjai laimi, jei žaidimas baigiasi per 20 raundus, kur vienas raundas apibrėžiamas kaip du žaidėjai, kiekvienas išsiuntęs po vieną galiojančią komandą.} \\
\\ 
\texttt{**Taškų skaičiavimas:**} \\
\\ 
\texttt{Kuo arčiau vienodų objektų abiejose žaidimo lentose, tuo daugiau taškų gausite. Baudos mažina jūsų taškus. Ar pavyks pagerinti rekordą?} \\
\\ 
\texttt{Kitas žaidėjas pradėjo žaidimą išsiųsdamas šią žinutę:} \\
\\ 
\texttt{"norėčiau susitarti dėl galutinės konfigūracijos. Prašau parodyti savo viziją, kur būtų išdėstyti objektai W, I, T, C, H."} \\
\\ 
\texttt{Kokia yra jūsų pirmoji komanda?} \\
            }
        }
    }
    & \\ \\

    \theutterance \stepcounter{utterance}  
    & & & \multicolumn{4}{p{0.6\linewidth}}{
        \cellcolor[rgb]{0.9,0.9,0.9}{
            \makecell[{{p{\linewidth}}}]{
                \texttt{\tiny{[GM$\langle$P2]}}
                \texttt{SAKYTI: Aš dar neturiu galutinės vizijos. Pirmiausia reikia susitarti dėl tikslinės konfigūracijos. Siųskiu savo lentos nuotrauką (naudoju tekstinį formatą): objektai yra išdėstyti taip: W – 6, I – 5, T – 1, C – 7, H – 2. Ar ši konfigūracija jums tinka kaip bendras tikslas?} \\
            }
        }
    }
     \\ \\

    \theutterance \stepcounter{utterance}  
    & & & \multicolumn{2}{p{0.3\linewidth}}{
        \cellcolor[rgb]{0.9,0.9,0.9}{
            \makecell[{{p{\linewidth}}}]{
                \texttt{\tiny{[GM$|$GM]}}
                \texttt{SAKYTI: Aš dar neturiu galutinės vizijos. Pirmiausia reikia susitarti dėl tikslinės konfigūracijos. Siųskiu savo lentos nuotrauką (naudoju tekstinį formatą): objektai yra išdėstyti taip: W – 6, I – 5, T – 1, C – 7, H – 2. Ar ši konfigūracija jums tinka kaip bendras tikslas?} \\
            }
        }
    }
    & & \\ \\

    \theutterance \stepcounter{utterance}  
    & & \multicolumn{4}{p{0.6\linewidth}}{
        \cellcolor[rgb]{0.9,0.9,0.9}{
            \makecell[{{p{\linewidth}}}]{
                \texttt{\tiny{[P1$\langle$GM]}}
                \texttt{Jūsų žinutė buvo perduota kitam žaidėjui.} \\
\\ 
\\ 
\texttt{Šiuo metu žaidžiate 1 raundą iš maksimaliai 20.} \\
\texttt{Jūs bendrai sukaupėte 0 iš 12 baudų.} \\
\texttt{Kitas žaidėjas atsiuntė šią žinutę:} \\
\\ 
\\ 
\texttt{"Aš dar neturiu galutinės vizijos. Pirmiausia reikia susitarti dėl tikslinės konfigūracijos. Siųskiu savo lentos nuotrauką (naudoju tekstinį formatą): objektai yra išdėstyti taip: W – 6, I – 5, T – 1, C – 7, H – 2. Ar ši konfigūracija jums tinka kaip bendras tikslas?"} \\
\\ 
\\ 
\texttt{Kokia jūsų kita komanda?} \\
            }
        }
    }
    & \\ \\

    \theutterance \stepcounter{utterance}  
    & \multicolumn{4}{p{0.6\linewidth}}{
        \cellcolor[rgb]{0.9,0.9,0.9}{
            \makecell[{{p{\linewidth}}}]{
                \texttt{\tiny{[P1$\rangle$GM]}}
                \texttt{SAKYTI: Taip, ši konfigūracija mane tenkina kaip bendras tikslas. Tad planuoju perkelti objektus taip: W į 6, I į 5, T į 1, C į 7, H į 2. Pradėsiu su objektu T, perkeldamas jį į langelį 1, 4.} \\
            }
        }
    }
    & & \\ \\

    \theutterance \stepcounter{utterance}  
    & & & \multicolumn{2}{p{0.3\linewidth}}{
        \cellcolor[rgb]{0.9,0.9,0.9}{
            \makecell[{{p{\linewidth}}}]{
                \texttt{\tiny{[GM$|$GM]}}
                \texttt{SAKYTI: Taip, ši konfigūracija mane tenkina kaip bendras tikslas. Tad planuoju perkelti objektus taip: W į 6, I į 5, T į 1, C į 7, H į 2. Pradėsiu su objektu T, perkeldamas jį į langelį 1, 4.} \\
            }
        }
    }
    & & \\ \\

    \theutterance \stepcounter{utterance}  
    & & \multicolumn{4}{p{0.6\linewidth}}{
        \cellcolor[rgb]{0.9,0.9,0.9}{
            \makecell[{{p{\linewidth}}}]{
                \texttt{\tiny{[GM$\rangle$P2]}}
                \texttt{Jūsų žinutė buvo perduota kitam žaidėjui.} \\
\\ 
\\ 
\texttt{Šiuo metu žaidžiate 2 raundą iš maksimaliai 20.} \\
\texttt{Jūs bendrai sukaupėte 0 iš 12 baudų.} \\
\texttt{Kitas žaidėjas atsiuntė šią žinutę:} \\
\\ 
\\ 
\texttt{"Taip, ši konfigūracija mane tenkina kaip bendras tikslas. Tad planuoju perkelti objektus taip: W į 6, I į 5, T į 1, C į 7, H į 2. Pradėsiu su objektu T, perkeldamas jį į langelį 1, 4."} \\
\\ 
\\ 
\texttt{Kokia jūsų kita komanda?} \\
            }
        }
    }
    & \\ \\

    \theutterance \stepcounter{utterance}  
    & & & \multicolumn{4}{p{0.6\linewidth}}{
        \cellcolor[rgb]{0.9,0.9,0.9}{
            \makecell[{{p{\linewidth}}}]{
                \texttt{\tiny{[GM$\langle$P2]}}
                \texttt{JUDĖTI: T, (1, 4)} \\
            }
        }
    }
     \\ \\

    \theutterance \stepcounter{utterance}  
    & & & \multicolumn{2}{p{0.3\linewidth}}{
        \cellcolor[rgb]{0.9,0.9,0.9}{
            \makecell[{{p{\linewidth}}}]{
                \texttt{\tiny{[GM$|$GM]}}
                \texttt{JUDĖTI: T, (1, 4)} \\
            }
        }
    }
    & & \\ \\

    \theutterance \stepcounter{utterance}  
    & & & \multicolumn{2}{p{0.3\linewidth}}{
        \cellcolor[rgb]{0.9,0.9,0.9}{
            \makecell[{{p{\linewidth}}}]{
                \texttt{\tiny{[GM$|$GM]}}
                \texttt{Nuobauda: (1,4) nėra tuščias, bet jame yra '─'.} \\
\\ 
\texttt{Jūs bendrai sukaupėte 1 iš 12 baudų.} \\
\texttt{Pabandykite dar kartą!} \\
            }
        }
    }
    & & \\ \\

    \theutterance \stepcounter{utterance}  
    & & \multicolumn{4}{p{0.6\linewidth}}{
        \cellcolor[rgb]{0.9,0.9,0.9}{
            \makecell[{{p{\linewidth}}}]{
                \texttt{\tiny{[GM$\rangle$P2]}}
                \texttt{Nuobauda: (1,4) nėra tuščias, bet jame yra '─'.} \\
\\ 
\texttt{Jūs bendrai sukaupėte 1 iš 12 baudų.} \\
\texttt{Pabandykite dar kartą!} \\
            }
        }
    }
    & \\ \\

    \theutterance \stepcounter{utterance}  
    & & & \multicolumn{4}{p{0.6\linewidth}}{
        \cellcolor[rgb]{0.9,0.9,0.9}{
            \makecell[{{p{\linewidth}}}]{
                \texttt{\tiny{[GM$\langle$P2]}}
                \texttt{JUDĖTI: T, (1, 6)} \\
            }
        }
    }
     \\ \\

    \theutterance \stepcounter{utterance}  
    & & & \multicolumn{2}{p{0.3\linewidth}}{
        \cellcolor[rgb]{0.9,0.9,0.9}{
            \makecell[{{p{\linewidth}}}]{
                \texttt{\tiny{[GM$|$GM]}}
                \texttt{JUDĖTI: T, (1, 6)} \\
            }
        }
    }
    & & \\ \\

    \theutterance \stepcounter{utterance}  
    & & & \multicolumn{2}{p{0.3\linewidth}}{
        \cellcolor[rgb]{0.9,0.9,0.9}{
            \makecell[{{p{\linewidth}}}]{
                \texttt{\tiny{[GM$|$GM]}}
                \texttt{T moved to (1, 6)} \\
            }
        }
    }
    & & \\ \\

    \theutterance \stepcounter{utterance}  
    & & \multicolumn{4}{p{0.6\linewidth}}{
        \cellcolor[rgb]{0.9,0.9,0.9}{
            \makecell[{{p{\linewidth}}}]{
                \texttt{\tiny{[P1$\langle$GM]}}
                \texttt{Jūsų žinutė buvo perduota kitam žaidėjui.} \\
\\ 
\\ 
\texttt{Šiuo metu žaidžiate 2 raundą iš maksimaliai 20.} \\
\texttt{Jūs bendrai sukaupėte 1 iš 12 baudų.} \\
\texttt{Kitas žaidėjas perkėlė objektą savo tinklelyje.} \\
\\ 
\\ 
\texttt{Kokia jūsų kita komanda?} \\
            }
        }
    }
    & \\ \\

    \theutterance \stepcounter{utterance}  
    & \multicolumn{4}{p{0.6\linewidth}}{
        \cellcolor[rgb]{0.9,0.9,0.9}{
            \makecell[{{p{\linewidth}}}]{
                \texttt{\tiny{[P1$\rangle$GM]}}
                \texttt{JUDĖTI: T, (1, 4)} \\
            }
        }
    }
    & & \\ \\

    \theutterance \stepcounter{utterance}  
    & & & \multicolumn{2}{p{0.3\linewidth}}{
        \cellcolor[rgb]{0.9,0.9,0.9}{
            \makecell[{{p{\linewidth}}}]{
                \texttt{\tiny{[GM$|$GM]}}
                \texttt{JUDĖTI: T, (1, 4)} \\
            }
        }
    }
    & & \\ \\

    \theutterance \stepcounter{utterance}  
    & & & \multicolumn{2}{p{0.3\linewidth}}{
        \cellcolor[rgb]{0.9,0.9,0.9}{
            \makecell[{{p{\linewidth}}}]{
                \texttt{\tiny{[GM$|$GM]}}
                \texttt{Nuobauda: (1,4) nėra tuščias, bet jame yra '─'.} \\
\\ 
\texttt{Jūs bendrai sukaupėte 2 iš 12 baudų.} \\
\texttt{Pabandykite dar kartą!} \\
            }
        }
    }
    & & \\ \\

    \theutterance \stepcounter{utterance}  
    & & \multicolumn{4}{p{0.6\linewidth}}{
        \cellcolor[rgb]{0.9,0.9,0.9}{
            \makecell[{{p{\linewidth}}}]{
                \texttt{\tiny{[P1$\langle$GM]}}
                \texttt{Nuobauda: (1,4) nėra tuščias, bet jame yra '─'.} \\
\\ 
\texttt{Jūs bendrai sukaupėte 2 iš 12 baudų.} \\
\texttt{Pabandykite dar kartą!} \\
            }
        }
    }
    & \\ \\

    \theutterance \stepcounter{utterance}  
    & \multicolumn{4}{p{0.6\linewidth}}{
        \cellcolor[rgb]{0.9,0.9,0.9}{
            \makecell[{{p{\linewidth}}}]{
                \texttt{\tiny{[P1$\rangle$GM]}}
                \texttt{JUDĖTI: T, (3, 4)} \\
            }
        }
    }
    & & \\ \\

    \theutterance \stepcounter{utterance}  
    & & & \multicolumn{2}{p{0.3\linewidth}}{
        \cellcolor[rgb]{0.9,0.9,0.9}{
            \makecell[{{p{\linewidth}}}]{
                \texttt{\tiny{[GM$|$GM]}}
                \texttt{JUDĖTI: T, (3, 4)} \\
            }
        }
    }
    & & \\ \\

    \theutterance \stepcounter{utterance}  
    & & & \multicolumn{2}{p{0.3\linewidth}}{
        \cellcolor[rgb]{0.9,0.9,0.9}{
            \makecell[{{p{\linewidth}}}]{
                \texttt{\tiny{[GM$|$GM]}}
                \texttt{Nuobauda: (3,4) nėra tuščias, bet jame yra '┐'.} \\
\\ 
\texttt{Jūs bendrai sukaupėte 3 iš 12 baudų.} \\
\texttt{Pabandykite dar kartą!} \\
            }
        }
    }
    & & \\ \\

    \theutterance \stepcounter{utterance}  
    & & \multicolumn{4}{p{0.6\linewidth}}{
        \cellcolor[rgb]{0.9,0.9,0.9}{
            \makecell[{{p{\linewidth}}}]{
                \texttt{\tiny{[P1$\langle$GM]}}
                \texttt{Nuobauda: (3,4) nėra tuščias, bet jame yra '┐'.} \\
\\ 
\texttt{Jūs bendrai sukaupėte 3 iš 12 baudų.} \\
\texttt{Pabandykite dar kartą!} \\
            }
        }
    }
    & \\ \\

    \theutterance \stepcounter{utterance}  
    & \multicolumn{4}{p{0.6\linewidth}}{
        \cellcolor[rgb]{0.9,0.9,0.9}{
            \makecell[{{p{\linewidth}}}]{
                \texttt{\tiny{[P1$\rangle$GM]}}
                \texttt{JUDĖTI: T, (5, 4)} \\
            }
        }
    }
    & & \\ \\

    \theutterance \stepcounter{utterance}  
    & & & \multicolumn{2}{p{0.3\linewidth}}{
        \cellcolor[rgb]{0.9,0.9,0.9}{
            \makecell[{{p{\linewidth}}}]{
                \texttt{\tiny{[GM$|$GM]}}
                \texttt{JUDĖTI: T, (5, 4)} \\
            }
        }
    }
    & & \\ \\

    \theutterance \stepcounter{utterance}  
    & & & \multicolumn{2}{p{0.3\linewidth}}{
        \cellcolor[rgb]{0.9,0.9,0.9}{
            \makecell[{{p{\linewidth}}}]{
                \texttt{\tiny{[GM$|$GM]}}
                \texttt{T moved to (5, 4)} \\
            }
        }
    }
    & & \\ \\

    \theutterance \stepcounter{utterance}  
    & & \multicolumn{4}{p{0.6\linewidth}}{
        \cellcolor[rgb]{0.9,0.9,0.9}{
            \makecell[{{p{\linewidth}}}]{
                \texttt{\tiny{[GM$\rangle$P2]}}
                \texttt{Sėkmingai perkelta 'T' į (1,6). Jūsų atnaujintas tinklelis atrodo taip:} \\
\texttt{```} \\
\texttt{ 1234567} \\
\texttt{╔═══════╗} \\
\texttt{║◌◌◌◌◌H◌║ 1} \\
\texttt{║◌◌◌◌◌◌◌║ 2} \\
\texttt{║◌◌◌◌◌┌─╢ 3} \\
\texttt{╟──┐◌◌│◌║ 4} \\
\texttt{║I◌├──┼─╢ 5} \\
\texttt{║TW│◌◌│C║ 6} \\
\texttt{║◌◌│◌◌│◌║ 7} \\
\texttt{╚══╧══╧═╝} \\
\texttt{```} \\
\texttt{Šiuo metu žaidžiate 3 raundą iš maksimaliai 20.} \\
\texttt{Jūs bendrai sukaupėte 3 iš 12 baudų.} \\
\texttt{Kitas žaidėjas perkėlė objektą savo tinklelyje.} \\
\\ 
\\ 
\texttt{Kokia jūsų kita komanda?} \\
            }
        }
    }
    & \\ \\

    \theutterance \stepcounter{utterance}  
    & & & \multicolumn{4}{p{0.6\linewidth}}{
        \cellcolor[rgb]{0.9,0.9,0.9}{
            \makecell[{{p{\linewidth}}}]{
                \texttt{\tiny{[GM$\langle$P2]}}
                \texttt{SAKYTI: Pradėsiu su objektu I, kuris dabar yra 5, planuoju perkelti jį į 3. Patikrinsiu, ar langelis (3,5) yra tuščias, ir jei taip – nurodysiu perkėlimą JUDĖTI: I, (3,5).} \\
            }
        }
    }
     \\ \\

    \theutterance \stepcounter{utterance}  
    & & & \multicolumn{2}{p{0.3\linewidth}}{
        \cellcolor[rgb]{0.9,0.9,0.9}{
            \makecell[{{p{\linewidth}}}]{
                \texttt{\tiny{[GM$|$GM]}}
                \texttt{SAKYTI: Pradėsiu su objektu I, kuris dabar yra 5, planuoju perkelti jį į 3. Patikrinsiu, ar langelis (3,5) yra tuščias, ir jei taip – nurodysiu perkėlimą JUDĖTI: I, (3,5).} \\
            }
        }
    }
    & & \\ \\

    \theutterance \stepcounter{utterance}  
    & & & \multicolumn{2}{p{0.3\linewidth}}{
        \cellcolor[rgb]{0.9,0.9,0.9}{
            \makecell[{{p{\linewidth}}}]{
                \texttt{\tiny{[GM$|$GM]}}
                \texttt{Invalid response format: SAKYTI: Pradėsiu su objektu I, kuris dabar yra 5, planuoju perkelti jį į 3. Patikrinsiu, ar langelis (3,5) yra tuščias, ir jei taip – nurodysiu perkėlimą JUDĖTI: I, (3,5).} \\
            }
        }
    }
    & & \\ \\

    \theutterance \stepcounter{utterance}  
    & & \multicolumn{4}{p{0.6\linewidth}}{
        \cellcolor[rgb]{0.9,0.9,0.9}{
            \makecell[{{p{\linewidth}}}]{
                \texttt{\tiny{[GM$\rangle$P2]}}
                \texttt{Nuobauda: Jūsų žinutėje yra daugiau nei viena komanda!} \\
\\ 
\texttt{Įsitikinkite, kad jūsų atsakyme yra tik `SAKYTI: <ŽINUTĖ>` arba `JUDĖTI: <OBJEKTAS>, (<X>, <Y>)`, ir nieko daugiau!} \\
\\ 
\texttt{Jūs bendrai sukaupėte 4 iš 12 baudų.} \\
\texttt{Pabandykite dar kartą!} \\
            }
        }
    }
    & \\ \\

    \theutterance \stepcounter{utterance}  
    & & & \multicolumn{4}{p{0.6\linewidth}}{
        \cellcolor[rgb]{0.9,0.9,0.9}{
            \makecell[{{p{\linewidth}}}]{
                \texttt{\tiny{[GM$\langle$P2]}}
                \texttt{SAKYTI: Pradėsiu su objektu I, planuoju perkelti į (3,5). Patikrinsiu langelio tuštumą prieš perdavimą.} \\
            }
        }
    }
     \\ \\

    \theutterance \stepcounter{utterance}  
    & & & \multicolumn{2}{p{0.3\linewidth}}{
        \cellcolor[rgb]{0.9,0.9,0.9}{
            \makecell[{{p{\linewidth}}}]{
                \texttt{\tiny{[GM$|$GM]}}
                \texttt{SAKYTI: Pradėsiu su objektu I, planuoju perkelti į (3,5). Patikrinsiu langelio tuštumą prieš perdavimą.} \\
            }
        }
    }
    & & \\ \\

    \theutterance \stepcounter{utterance}  
    & & \multicolumn{4}{p{0.6\linewidth}}{
        \cellcolor[rgb]{0.9,0.9,0.9}{
            \makecell[{{p{\linewidth}}}]{
                \texttt{\tiny{[P1$\langle$GM]}}
                \texttt{Sėkmingai perkelta 'T' į (5,4). Jūsų atnaujintas tinklelis atrodo taip:} \\
\texttt{```} \\
\texttt{ 1234567} \\
\texttt{╔═══════╗} \\
\texttt{║◌◌◌◌◌◌◌║ 1} \\
\texttt{║◌◌◌◌◌◌◌║ 2} \\
\texttt{║◌◌W◌◌┌─╢ 3} \\
\texttt{╟──┐◌T│H║ 4} \\
\texttt{║C◌├──┼─╢ 5} \\
\texttt{║◌◌│I◌│◌║ 6} \\
\texttt{║◌◌│◌◌│◌║ 7} \\
\texttt{╚══╧══╧═╝} \\
\texttt{```} \\
\texttt{Šiuo metu žaidžiate 3 raundą iš maksimaliai 20.} \\
\texttt{Jūs bendrai sukaupėte 4 iš 12 baudų.} \\
\texttt{Kitas žaidėjas atsiuntė šią žinutę:} \\
\\ 
\\ 
\texttt{"Pradėsiu su objektu I, planuoju perkelti į (3,5). Patikrinsiu langelio tuštumą prieš perdavimą."} \\
\\ 
\\ 
\texttt{Kokia jūsų kita komanda?} \\
            }
        }
    }
    & \\ \\

    \theutterance \stepcounter{utterance}  
    & \multicolumn{4}{p{0.6\linewidth}}{
        \cellcolor[rgb]{0.9,0.9,0.9}{
            \makecell[{{p{\linewidth}}}]{
                \texttt{\tiny{[P1$\rangle$GM]}}
                \texttt{SAKYTI: Patvirtinu objektą I perkeltą į (3,5). Ačiū už informaciją.} \\
            }
        }
    }
    & & \\ \\

    \theutterance \stepcounter{utterance}  
    & & & \multicolumn{2}{p{0.3\linewidth}}{
        \cellcolor[rgb]{0.9,0.9,0.9}{
            \makecell[{{p{\linewidth}}}]{
                \texttt{\tiny{[GM$|$GM]}}
                \texttt{SAKYTI: Patvirtinu objektą I perkeltą į (3,5). Ačiū už informaciją.} \\
            }
        }
    }
    & & \\ \\

    \theutterance \stepcounter{utterance}  
    & & \multicolumn{4}{p{0.6\linewidth}}{
        \cellcolor[rgb]{0.9,0.9,0.9}{
            \makecell[{{p{\linewidth}}}]{
                \texttt{\tiny{[GM$\rangle$P2]}}
                \texttt{Jūsų žinutė buvo perduota kitam žaidėjui.} \\
\\ 
\\ 
\texttt{Šiuo metu žaidžiate 4 raundą iš maksimaliai 20.} \\
\texttt{Jūs bendrai sukaupėte 4 iš 12 baudų.} \\
\texttt{Kitas žaidėjas atsiuntė šią žinutę:} \\
\\ 
\\ 
\texttt{"Patvirtinu objektą I perkeltą į (3,5). Ačiū už informaciją."} \\
\\ 
\\ 
\texttt{Kokia jūsų kita komanda?} \\
            }
        }
    }
    & \\ \\

    \theutterance \stepcounter{utterance}  
    & & & \multicolumn{4}{p{0.6\linewidth}}{
        \cellcolor[rgb]{0.9,0.9,0.9}{
            \makecell[{{p{\linewidth}}}]{
                \texttt{\tiny{[GM$\langle$P2]}}
                \texttt{JUDĖTI: I, (3,5)} \\
            }
        }
    }
     \\ \\

    \theutterance \stepcounter{utterance}  
    & & & \multicolumn{2}{p{0.3\linewidth}}{
        \cellcolor[rgb]{0.9,0.9,0.9}{
            \makecell[{{p{\linewidth}}}]{
                \texttt{\tiny{[GM$|$GM]}}
                \texttt{JUDĖTI: I, (3,5)} \\
            }
        }
    }
    & & \\ \\

    \theutterance \stepcounter{utterance}  
    & & & \multicolumn{2}{p{0.3\linewidth}}{
        \cellcolor[rgb]{0.9,0.9,0.9}{
            \makecell[{{p{\linewidth}}}]{
                \texttt{\tiny{[GM$|$GM]}}
                \texttt{Nuobauda: (3,5) nėra tuščias, bet jame yra '├'.} \\
\\ 
\texttt{Jūs bendrai sukaupėte 5 iš 12 baudų.} \\
\texttt{Pabandykite dar kartą!} \\
            }
        }
    }
    & & \\ \\

    \theutterance \stepcounter{utterance}  
    & & \multicolumn{4}{p{0.6\linewidth}}{
        \cellcolor[rgb]{0.9,0.9,0.9}{
            \makecell[{{p{\linewidth}}}]{
                \texttt{\tiny{[GM$\rangle$P2]}}
                \texttt{Nuobauda: (3,5) nėra tuščias, bet jame yra '├'.} \\
\\ 
\texttt{Jūs bendrai sukaupėte 5 iš 12 baudų.} \\
\texttt{Pabandykite dar kartą!} \\
            }
        }
    }
    & \\ \\

    \theutterance \stepcounter{utterance}  
    & & & \multicolumn{4}{p{0.6\linewidth}}{
        \cellcolor[rgb]{0.9,0.9,0.9}{
            \makecell[{{p{\linewidth}}}]{
                \texttt{\tiny{[GM$\langle$P2]}}
                \texttt{SAKYTI: Pradėsiu su objektu H, kuris dabar yra 2. Planuoju perkelti jį į 4. Patikrinsiu langelio (4,2) tuštumą prieš perdavimą.} \\
            }
        }
    }
     \\ \\

    \theutterance \stepcounter{utterance}  
    & & & \multicolumn{2}{p{0.3\linewidth}}{
        \cellcolor[rgb]{0.9,0.9,0.9}{
            \makecell[{{p{\linewidth}}}]{
                \texttt{\tiny{[GM$|$GM]}}
                \texttt{SAKYTI: Pradėsiu su objektu H, kuris dabar yra 2. Planuoju perkelti jį į 4. Patikrinsiu langelio (4,2) tuštumą prieš perdavimą.} \\
            }
        }
    }
    & & \\ \\

    \theutterance \stepcounter{utterance}  
    & & \multicolumn{4}{p{0.6\linewidth}}{
        \cellcolor[rgb]{0.9,0.9,0.9}{
            \makecell[{{p{\linewidth}}}]{
                \texttt{\tiny{[P1$\langle$GM]}}
                \texttt{Jūsų žinutė buvo perduota kitam žaidėjui.} \\
\\ 
\\ 
\texttt{Šiuo metu žaidžiate 4 raundą iš maksimaliai 20.} \\
\texttt{Jūs bendrai sukaupėte 5 iš 12 baudų.} \\
\texttt{Kitas žaidėjas atsiuntė šią žinutę:} \\
\\ 
\\ 
\texttt{"Pradėsiu su objektu H, kuris dabar yra 2. Planuoju perkelti jį į 4. Patikrinsiu langelio (4,2) tuštumą prieš perdavimą."} \\
\\ 
\\ 
\texttt{Kokia jūsų kita komanda?} \\
            }
        }
    }
    & \\ \\

    \theutterance \stepcounter{utterance}  
    & \multicolumn{4}{p{0.6\linewidth}}{
        \cellcolor[rgb]{0.9,0.9,0.9}{
            \makecell[{{p{\linewidth}}}]{
                \texttt{\tiny{[P1$\rangle$GM]}}
                \texttt{JUDĖTI: H, (4, 2)} \\
            }
        }
    }
    & & \\ \\

    \theutterance \stepcounter{utterance}  
    & & & \multicolumn{2}{p{0.3\linewidth}}{
        \cellcolor[rgb]{0.9,0.9,0.9}{
            \makecell[{{p{\linewidth}}}]{
                \texttt{\tiny{[GM$|$GM]}}
                \texttt{JUDĖTI: H, (4, 2)} \\
            }
        }
    }
    & & \\ \\

    \theutterance \stepcounter{utterance}  
    & & & \multicolumn{2}{p{0.3\linewidth}}{
        \cellcolor[rgb]{0.9,0.9,0.9}{
            \makecell[{{p{\linewidth}}}]{
                \texttt{\tiny{[GM$|$GM]}}
                \texttt{H moved to (4, 2)} \\
            }
        }
    }
    & & \\ \\

    \theutterance \stepcounter{utterance}  
    & & \multicolumn{4}{p{0.6\linewidth}}{
        \cellcolor[rgb]{0.9,0.9,0.9}{
            \makecell[{{p{\linewidth}}}]{
                \texttt{\tiny{[GM$\rangle$P2]}}
                \texttt{Jūsų žinutė buvo perduota kitam žaidėjui.} \\
\\ 
\\ 
\texttt{Šiuo metu žaidžiate 5 raundą iš maksimaliai 20.} \\
\texttt{Jūs bendrai sukaupėte 5 iš 12 baudų.} \\
\texttt{Kitas žaidėjas perkėlė objektą savo tinklelyje.} \\
\\ 
\\ 
\texttt{Kokia jūsų kita komanda?} \\
            }
        }
    }
    & \\ \\

    \theutterance \stepcounter{utterance}  
    & & & \multicolumn{4}{p{0.6\linewidth}}{
        \cellcolor[rgb]{0.9,0.9,0.9}{
            \makecell[{{p{\linewidth}}}]{
                \texttt{\tiny{[GM$\langle$P2]}}
                \texttt{SAKYTI: Pradėsiu su objektu T, kuris dabar yra 6. Planuoju perkelti jį į 7, norėdamas užpildyti langelį. Patikrinsiu, ar langelis (7,6) yra tuščias, kol perduosiu komandinę komandą.} \\
            }
        }
    }
     \\ \\

    \theutterance \stepcounter{utterance}  
    & & & \multicolumn{2}{p{0.3\linewidth}}{
        \cellcolor[rgb]{0.9,0.9,0.9}{
            \makecell[{{p{\linewidth}}}]{
                \texttt{\tiny{[GM$|$GM]}}
                \texttt{SAKYTI: Pradėsiu su objektu T, kuris dabar yra 6. Planuoju perkelti jį į 7, norėdamas užpildyti langelį. Patikrinsiu, ar langelis (7,6) yra tuščias, kol perduosiu komandinę komandą.} \\
            }
        }
    }
    & & \\ \\

    \theutterance \stepcounter{utterance}  
    & & \multicolumn{4}{p{0.6\linewidth}}{
        \cellcolor[rgb]{0.9,0.9,0.9}{
            \makecell[{{p{\linewidth}}}]{
                \texttt{\tiny{[P1$\langle$GM]}}
                \texttt{Sėkmingai perkelta 'H' į (4,2). Jūsų atnaujintas tinklelis atrodo taip:} \\
\texttt{```} \\
\texttt{ 1234567} \\
\texttt{╔═══════╗} \\
\texttt{║◌◌◌◌◌◌◌║ 1} \\
\texttt{║◌◌◌H◌◌◌║ 2} \\
\texttt{║◌◌W◌◌┌─╢ 3} \\
\texttt{╟──┐◌T│◌║ 4} \\
\texttt{║C◌├──┼─╢ 5} \\
\texttt{║◌◌│I◌│◌║ 6} \\
\texttt{║◌◌│◌◌│◌║ 7} \\
\texttt{╚══╧══╧═╝} \\
\texttt{```} \\
\texttt{Šiuo metu žaidžiate 5 raundą iš maksimaliai 20.} \\
\texttt{Jūs bendrai sukaupėte 5 iš 12 baudų.} \\
\texttt{Kitas žaidėjas atsiuntė šią žinutę:} \\
\\ 
\\ 
\texttt{"Pradėsiu su objektu T, kuris dabar yra 6. Planuoju perkelti jį į 7, norėdamas užpildyti langelį. Patikrinsiu, ar langelis (7,6) yra tuščias, kol perduosiu komandinę komandą."} \\
\\ 
\\ 
\texttt{Kokia jūsų kita komanda?} \\
            }
        }
    }
    & \\ \\

    \theutterance \stepcounter{utterance}  
    & \multicolumn{4}{p{0.6\linewidth}}{
        \cellcolor[rgb]{0.9,0.9,0.9}{
            \makecell[{{p{\linewidth}}}]{
                \texttt{\tiny{[P1$\rangle$GM]}}
                \texttt{SAKYTI: Patvirtinu perkėlimą į (7,6).} \\
            }
        }
    }
    & & \\ \\

    \theutterance \stepcounter{utterance}  
    & & & \multicolumn{2}{p{0.3\linewidth}}{
        \cellcolor[rgb]{0.9,0.9,0.9}{
            \makecell[{{p{\linewidth}}}]{
                \texttt{\tiny{[GM$|$GM]}}
                \texttt{SAKYTI: Patvirtinu perkėlimą į (7,6).} \\
            }
        }
    }
    & & \\ \\

    \theutterance \stepcounter{utterance}  
    & & \multicolumn{4}{p{0.6\linewidth}}{
        \cellcolor[rgb]{0.9,0.9,0.9}{
            \makecell[{{p{\linewidth}}}]{
                \texttt{\tiny{[GM$\rangle$P2]}}
                \texttt{Jūsų žinutė buvo perduota kitam žaidėjui.} \\
\\ 
\\ 
\texttt{Šiuo metu žaidžiate 6 raundą iš maksimaliai 20.} \\
\texttt{Jūs bendrai sukaupėte 5 iš 12 baudų.} \\
\texttt{Kitas žaidėjas atsiuntė šią žinutę:} \\
\\ 
\\ 
\texttt{"Patvirtinu perkėlimą į (7,6)."} \\
\\ 
\\ 
\texttt{Kokia jūsų kita komanda?} \\
            }
        }
    }
    & \\ \\

    \theutterance \stepcounter{utterance}  
    & & & \multicolumn{4}{p{0.6\linewidth}}{
        \cellcolor[rgb]{0.9,0.9,0.9}{
            \makecell[{{p{\linewidth}}}]{
                \texttt{\tiny{[GM$\langle$P2]}}
                \texttt{JUDĖTI: T, (7,6)} \\
            }
        }
    }
     \\ \\

    \theutterance \stepcounter{utterance}  
    & & & \multicolumn{2}{p{0.3\linewidth}}{
        \cellcolor[rgb]{0.9,0.9,0.9}{
            \makecell[{{p{\linewidth}}}]{
                \texttt{\tiny{[GM$|$GM]}}
                \texttt{JUDĖTI: T, (7,6)} \\
            }
        }
    }
    & & \\ \\

    \theutterance \stepcounter{utterance}  
    & & & \multicolumn{2}{p{0.3\linewidth}}{
        \cellcolor[rgb]{0.9,0.9,0.9}{
            \makecell[{{p{\linewidth}}}]{
                \texttt{\tiny{[GM$|$GM]}}
                \texttt{Nuobauda: (7,6) nėra tuščias, bet jame yra 'C'.} \\
\\ 
\texttt{Jūs bendrai sukaupėte 6 iš 12 baudų.} \\
\texttt{Pabandykite dar kartą!} \\
            }
        }
    }
    & & \\ \\

    \theutterance \stepcounter{utterance}  
    & & \multicolumn{4}{p{0.6\linewidth}}{
        \cellcolor[rgb]{0.9,0.9,0.9}{
            \makecell[{{p{\linewidth}}}]{
                \texttt{\tiny{[GM$\rangle$P2]}}
                \texttt{Nuobauda: (7,6) nėra tuščias, bet jame yra 'C'.} \\
\\ 
\texttt{Jūs bendrai sukaupėte 6 iš 12 baudų.} \\
\texttt{Pabandykite dar kartą!} \\
            }
        }
    }
    & \\ \\

    \theutterance \stepcounter{utterance}  
    & & & \multicolumn{4}{p{0.6\linewidth}}{
        \cellcolor[rgb]{0.9,0.9,0.9}{
            \makecell[{{p{\linewidth}}}]{
                \texttt{\tiny{[GM$\langle$P2]}}
                \texttt{SAKYTI: Pradėsiu su objektu C, kuris dabar yra 7. Planuoju perkelti jį į 6, siekdamas atitikti ketinamą konfigūraciją. Patikrinsiu langelio (6,7) tuštumą prieš komandinę komandą.} \\
            }
        }
    }
     \\ \\

    \theutterance \stepcounter{utterance}  
    & & & \multicolumn{2}{p{0.3\linewidth}}{
        \cellcolor[rgb]{0.9,0.9,0.9}{
            \makecell[{{p{\linewidth}}}]{
                \texttt{\tiny{[GM$|$GM]}}
                \texttt{SAKYTI: Pradėsiu su objektu C, kuris dabar yra 7. Planuoju perkelti jį į 6, siekdamas atitikti ketinamą konfigūraciją. Patikrinsiu langelio (6,7) tuštumą prieš komandinę komandą.} \\
            }
        }
    }
    & & \\ \\

    \theutterance \stepcounter{utterance}  
    & & \multicolumn{4}{p{0.6\linewidth}}{
        \cellcolor[rgb]{0.9,0.9,0.9}{
            \makecell[{{p{\linewidth}}}]{
                \texttt{\tiny{[P1$\langle$GM]}}
                \texttt{Jūsų žinutė buvo perduota kitam žaidėjui.} \\
\\ 
\\ 
\texttt{Šiuo metu žaidžiate 6 raundą iš maksimaliai 20.} \\
\texttt{Jūs bendrai sukaupėte 6 iš 12 baudų.} \\
\texttt{Kitas žaidėjas atsiuntė šią žinutę:} \\
\\ 
\\ 
\texttt{"Pradėsiu su objektu C, kuris dabar yra 7. Planuoju perkelti jį į 6, siekdamas atitikti ketinamą konfigūraciją. Patikrinsiu langelio (6,7) tuštumą prieš komandinę komandą."} \\
\\ 
\\ 
\texttt{Kokia jūsų kita komanda?} \\
            }
        }
    }
    & \\ \\

    \theutterance \stepcounter{utterance}  
    & \multicolumn{4}{p{0.6\linewidth}}{
        \cellcolor[rgb]{0.9,0.9,0.9}{
            \makecell[{{p{\linewidth}}}]{
                \texttt{\tiny{[P1$\rangle$GM]}}
                \texttt{JUDĖTI: C, (6, 7)} \\
            }
        }
    }
    & & \\ \\

    \theutterance \stepcounter{utterance}  
    & & & \multicolumn{2}{p{0.3\linewidth}}{
        \cellcolor[rgb]{0.9,0.9,0.9}{
            \makecell[{{p{\linewidth}}}]{
                \texttt{\tiny{[GM$|$GM]}}
                \texttt{JUDĖTI: C, (6, 7)} \\
            }
        }
    }
    & & \\ \\

    \theutterance \stepcounter{utterance}  
    & & & \multicolumn{2}{p{0.3\linewidth}}{
        \cellcolor[rgb]{0.9,0.9,0.9}{
            \makecell[{{p{\linewidth}}}]{
                \texttt{\tiny{[GM$|$GM]}}
                \texttt{Nuobauda: (6,7) nėra tuščias, bet jame yra '│'.} \\
\\ 
\texttt{Jūs bendrai sukaupėte 7 iš 12 baudų.} \\
\texttt{Pabandykite dar kartą!} \\
            }
        }
    }
    & & \\ \\

    \theutterance \stepcounter{utterance}  
    & & \multicolumn{4}{p{0.6\linewidth}}{
        \cellcolor[rgb]{0.9,0.9,0.9}{
            \makecell[{{p{\linewidth}}}]{
                \texttt{\tiny{[P1$\langle$GM]}}
                \texttt{Nuobauda: (6,7) nėra tuščias, bet jame yra '│'.} \\
\\ 
\texttt{Jūs bendrai sukaupėte 7 iš 12 baudų.} \\
\texttt{Pabandykite dar kartą!} \\
            }
        }
    }
    & \\ \\

    \theutterance \stepcounter{utterance}  
    & \multicolumn{4}{p{0.6\linewidth}}{
        \cellcolor[rgb]{0.9,0.9,0.9}{
            \makecell[{{p{\linewidth}}}]{
                \texttt{\tiny{[P1$\rangle$GM]}}
                \texttt{JUDĖTI: C, (6, 5)} \\
            }
        }
    }
    & & \\ \\

    \theutterance \stepcounter{utterance}  
    & & & \multicolumn{2}{p{0.3\linewidth}}{
        \cellcolor[rgb]{0.9,0.9,0.9}{
            \makecell[{{p{\linewidth}}}]{
                \texttt{\tiny{[GM$|$GM]}}
                \texttt{JUDĖTI: C, (6, 5)} \\
            }
        }
    }
    & & \\ \\

    \theutterance \stepcounter{utterance}  
    & & & \multicolumn{2}{p{0.3\linewidth}}{
        \cellcolor[rgb]{0.9,0.9,0.9}{
            \makecell[{{p{\linewidth}}}]{
                \texttt{\tiny{[GM$|$GM]}}
                \texttt{Nuobauda: (6,5) nėra tuščias, bet jame yra '┼'.} \\
\\ 
\texttt{Jūs bendrai sukaupėte 8 iš 12 baudų.} \\
\texttt{Pabandykite dar kartą!} \\
            }
        }
    }
    & & \\ \\

    \theutterance \stepcounter{utterance}  
    & & \multicolumn{4}{p{0.6\linewidth}}{
        \cellcolor[rgb]{0.9,0.9,0.9}{
            \makecell[{{p{\linewidth}}}]{
                \texttt{\tiny{[P1$\langle$GM]}}
                \texttt{Nuobauda: (6,5) nėra tuščias, bet jame yra '┼'.} \\
\\ 
\texttt{Jūs bendrai sukaupėte 8 iš 12 baudų.} \\
\texttt{Pabandykite dar kartą!} \\
            }
        }
    }
    & \\ \\

    \theutterance \stepcounter{utterance}  
    & \multicolumn{4}{p{0.6\linewidth}}{
        \cellcolor[rgb]{0.9,0.9,0.9}{
            \makecell[{{p{\linewidth}}}]{
                \texttt{\tiny{[P1$\rangle$GM]}}
                \texttt{JUDĖTI: C, (5, 5)} \\
            }
        }
    }
    & & \\ \\

    \theutterance \stepcounter{utterance}  
    & & & \multicolumn{2}{p{0.3\linewidth}}{
        \cellcolor[rgb]{0.9,0.9,0.9}{
            \makecell[{{p{\linewidth}}}]{
                \texttt{\tiny{[GM$|$GM]}}
                \texttt{JUDĖTI: C, (5, 5)} \\
            }
        }
    }
    & & \\ \\

    \theutterance \stepcounter{utterance}  
    & & & \multicolumn{2}{p{0.3\linewidth}}{
        \cellcolor[rgb]{0.9,0.9,0.9}{
            \makecell[{{p{\linewidth}}}]{
                \texttt{\tiny{[GM$|$GM]}}
                \texttt{Nuobauda: (5,5) nėra tuščias, bet jame yra '─'.} \\
\\ 
\texttt{Jūs bendrai sukaupėte 9 iš 12 baudų.} \\
\texttt{Pabandykite dar kartą!} \\
            }
        }
    }
    & & \\ \\

    \theutterance \stepcounter{utterance}  
    & & \multicolumn{4}{p{0.6\linewidth}}{
        \cellcolor[rgb]{0.9,0.9,0.9}{
            \makecell[{{p{\linewidth}}}]{
                \texttt{\tiny{[P1$\langle$GM]}}
                \texttt{Nuobauda: (5,5) nėra tuščias, bet jame yra '─'.} \\
\\ 
\texttt{Jūs bendrai sukaupėte 9 iš 12 baudų.} \\
\texttt{Pabandykite dar kartą!} \\
            }
        }
    }
    & \\ \\

    \theutterance \stepcounter{utterance}  
    & \multicolumn{4}{p{0.6\linewidth}}{
        \cellcolor[rgb]{0.9,0.9,0.9}{
            \makecell[{{p{\linewidth}}}]{
                \texttt{\tiny{[P1$\rangle$GM]}}
                \texttt{JUDĖTI: C, (6, 5)} \\
            }
        }
    }
    & & \\ \\

    \theutterance \stepcounter{utterance}  
    & & & \multicolumn{2}{p{0.3\linewidth}}{
        \cellcolor[rgb]{0.9,0.9,0.9}{
            \makecell[{{p{\linewidth}}}]{
                \texttt{\tiny{[GM$|$GM]}}
                \texttt{JUDĖTI: C, (6, 5)} \\
            }
        }
    }
    & & \\ \\

    \theutterance \stepcounter{utterance}  
    & & & \multicolumn{2}{p{0.3\linewidth}}{
        \cellcolor[rgb]{0.9,0.9,0.9}{
            \makecell[{{p{\linewidth}}}]{
                \texttt{\tiny{[GM$|$GM]}}
                \texttt{Nuobauda: (6,5) nėra tuščias, bet jame yra '┼'.} \\
\\ 
\texttt{Jūs bendrai sukaupėte 10 iš 12 baudų.} \\
\texttt{Pabandykite dar kartą!} \\
            }
        }
    }
    & & \\ \\

    \theutterance \stepcounter{utterance}  
    & & \multicolumn{4}{p{0.6\linewidth}}{
        \cellcolor[rgb]{0.9,0.9,0.9}{
            \makecell[{{p{\linewidth}}}]{
                \texttt{\tiny{[P1$\langle$GM]}}
                \texttt{Nuobauda: (6,5) nėra tuščias, bet jame yra '┼'.} \\
\\ 
\texttt{Jūs bendrai sukaupėte 10 iš 12 baudų.} \\
\texttt{Pabandykite dar kartą!} \\
            }
        }
    }
    & \\ \\

    \theutterance \stepcounter{utterance}  
    & \multicolumn{4}{p{0.6\linewidth}}{
        \cellcolor[rgb]{0.9,0.9,0.9}{
            \makecell[{{p{\linewidth}}}]{
                \texttt{\tiny{[P1$\rangle$GM]}}
                \texttt{JUDĖTI: C, (5, 5)} \\
            }
        }
    }
    & & \\ \\

    \theutterance \stepcounter{utterance}  
    & & & \multicolumn{2}{p{0.3\linewidth}}{
        \cellcolor[rgb]{0.9,0.9,0.9}{
            \makecell[{{p{\linewidth}}}]{
                \texttt{\tiny{[GM$|$GM]}}
                \texttt{JUDĖTI: C, (5, 5)} \\
            }
        }
    }
    & & \\ \\

    \theutterance \stepcounter{utterance}  
    & & & \multicolumn{2}{p{0.3\linewidth}}{
        \cellcolor[rgb]{0.9,0.9,0.9}{
            \makecell[{{p{\linewidth}}}]{
                \texttt{\tiny{[GM$|$GM]}}
                \texttt{Nuobauda: (5,5) nėra tuščias, bet jame yra '─'.} \\
\\ 
\texttt{Jūs bendrai sukaupėte 11 iš 12 baudų.} \\
\texttt{Pabandykite dar kartą!} \\
            }
        }
    }
    & & \\ \\

    \theutterance \stepcounter{utterance}  
    & & \multicolumn{4}{p{0.6\linewidth}}{
        \cellcolor[rgb]{0.9,0.9,0.9}{
            \makecell[{{p{\linewidth}}}]{
                \texttt{\tiny{[P1$\langle$GM]}}
                \texttt{Nuobauda: (5,5) nėra tuščias, bet jame yra '─'.} \\
\\ 
\texttt{Jūs bendrai sukaupėte 11 iš 12 baudų.} \\
\texttt{Pabandykite dar kartą!} \\
            }
        }
    }
    & \\ \\

    \theutterance \stepcounter{utterance}  
    & \multicolumn{4}{p{0.6\linewidth}}{
        \cellcolor[rgb]{0.9,0.9,0.9}{
            \makecell[{{p{\linewidth}}}]{
                \texttt{\tiny{[P1$\rangle$GM]}}
                \texttt{JUDĖTI: C, (4, 5)} \\
            }
        }
    }
    & & \\ \\

    \theutterance \stepcounter{utterance}  
    & & & \multicolumn{2}{p{0.3\linewidth}}{
        \cellcolor[rgb]{0.9,0.9,0.9}{
            \makecell[{{p{\linewidth}}}]{
                \texttt{\tiny{[GM$|$GM]}}
                \texttt{JUDĖTI: C, (4, 5)} \\
            }
        }
    }
    & & \\ \\

    \theutterance \stepcounter{utterance}  
    & & & \multicolumn{2}{p{0.3\linewidth}}{
        \cellcolor[rgb]{0.9,0.9,0.9}{
            \makecell[{{p{\linewidth}}}]{
                \texttt{\tiny{[GM$|$GM]}}
                \texttt{Nuobauda: (4,5) nėra tuščias, bet jame yra '─'.} \\
\\ 
\texttt{Jūs bendrai sukaupėte 12 iš 12 baudų.} \\
\texttt{Pabandykite dar kartą!} \\
            }
        }
    }
    & & \\ \\

    \theutterance \stepcounter{utterance}  
    & & \multicolumn{4}{p{0.6\linewidth}}{
        \cellcolor[rgb]{0.9,0.9,0.9}{
            \makecell[{{p{\linewidth}}}]{
                \texttt{\tiny{[P1$\langle$GM]}}
                \texttt{Nuobauda: (4,5) nėra tuščias, bet jame yra '─'.} \\
\\ 
\texttt{Jūs bendrai sukaupėte 12 iš 12 baudų.} \\
\texttt{Pabandykite dar kartą!} \\
            }
        }
    }
    & \\ \\

    \theutterance \stepcounter{utterance}  
    & \multicolumn{4}{p{0.6\linewidth}}{
        \cellcolor[rgb]{0.9,0.9,0.9}{
            \makecell[{{p{\linewidth}}}]{
                \texttt{\tiny{[P1$\rangle$GM]}}
                \texttt{SAKYTI: Atsiprašau, bet baudos riba pasiektą. Turėsiu lyginti šiuo metu turimas pozicijas su kitų žaidėjo lentos pozicijomis ir pateikti galutinę komandą.} \\
            }
        }
    }
    & & \\ \\

    \theutterance \stepcounter{utterance}  
    & & & \multicolumn{2}{p{0.3\linewidth}}{
        \cellcolor[rgb]{0.9,0.9,0.9}{
            \makecell[{{p{\linewidth}}}]{
                \texttt{\tiny{[GM$|$GM]}}
                \texttt{SAKYTI: Atsiprašau, bet baudos riba pasiektą. Turėsiu lyginti šiuo metu turimas pozicijas su kitų žaidėjo lentos pozicijomis ir pateikti galutinę komandą.} \\
            }
        }
    }
    & & \\ \\

    \theutterance \stepcounter{utterance}  
    & & \multicolumn{4}{p{0.6\linewidth}}{
        \cellcolor[rgb]{0.9,0.9,0.9}{
            \makecell[{{p{\linewidth}}}]{
                \texttt{\tiny{[GM$\rangle$P2]}}
                \texttt{Jūsų žinutė buvo perduota kitam žaidėjui.} \\
\\ 
\\ 
\texttt{Šiuo metu žaidžiate 7 raundą iš maksimaliai 20.} \\
\texttt{Jūs bendrai sukaupėte 12 iš 12 baudų.} \\
\texttt{Kitas žaidėjas atsiuntė šią žinutę:} \\
\\ 
\\ 
\texttt{"Atsiprašau, bet baudos riba pasiektą. Turėsiu lyginti šiuo metu turimas pozicijas su kitų žaidėjo lentos pozicijomis ir pateikti galutinę komandą."} \\
\\ 
\\ 
\texttt{Kokia jūsų kita komanda?} \\
            }
        }
    }
    & \\ \\

    \theutterance \stepcounter{utterance}  
    & & & \multicolumn{4}{p{0.6\linewidth}}{
        \cellcolor[rgb]{0.9,0.9,0.9}{
            \makecell[{{p{\linewidth}}}]{
                \texttt{\tiny{[GM$\langle$P2]}}
                \texttt{SAKYTI: Aš taip pat pasiekiau baudų ribą ir negaliu rizikuoti daugiau bandymų. Pradėsiu nuo objekto I, esančio 5, ir bandysiu perkelti į 3. Jei langelis (3,5) bus tuščias – pasiųsiu JUDĖTI: I, (3,5). Prieš tai dar kartą patikrinsiu langelio būseną.} \\
            }
        }
    }
     \\ \\

    \theutterance \stepcounter{utterance}  
    & & & \multicolumn{2}{p{0.3\linewidth}}{
        \cellcolor[rgb]{0.9,0.9,0.9}{
            \makecell[{{p{\linewidth}}}]{
                \texttt{\tiny{[GM$|$GM]}}
                \texttt{SAKYTI: Aš taip pat pasiekiau baudų ribą ir negaliu rizikuoti daugiau bandymų. Pradėsiu nuo objekto I, esančio 5, ir bandysiu perkelti į 3. Jei langelis (3,5) bus tuščias – pasiųsiu JUDĖTI: I, (3,5). Prieš tai dar kartą patikrinsiu langelio būseną.} \\
            }
        }
    }
    & & \\ \\

    \theutterance \stepcounter{utterance}  
    & & & \multicolumn{2}{p{0.3\linewidth}}{
        \cellcolor[rgb]{0.9,0.9,0.9}{
            \makecell[{{p{\linewidth}}}]{
                \texttt{\tiny{[GM$|$GM]}}
                \texttt{Invalid response format: SAKYTI: Aš taip pat pasiekiau baudų ribą ir negaliu rizikuoti daugiau bandymų. Pradėsiu nuo objekto I, esančio 5, ir bandysiu perkelti į 3. Jei langelis (3,5) bus tuščias – pasiųsiu JUDĖTI: I, (3,5). Prieš tai dar kartą patikrinsiu langelio būseną.} \\
            }
        }
    }
    & & \\ \\

    \theutterance \stepcounter{utterance}  
    & & & \multicolumn{2}{p{0.3\linewidth}}{
        \cellcolor[rgb]{0.9,0.9,0.9}{
            \makecell[{{p{\linewidth}}}]{
                \texttt{\tiny{[GM$|$GM]}}
                \texttt{Maximum number of penalties exceeded} \\
            }
        }
    }
    & & \\ \\

    \theutterance \stepcounter{utterance}  
    & & & \multicolumn{2}{p{0.3\linewidth}}{
        \cellcolor[rgb]{0.9,0.9,0.9}{
            \makecell[{{p{\linewidth}}}]{
                \texttt{\tiny{[GM$|$GM]}}
                \texttt{Initial states:} \\
\texttt{```} \\
\texttt{Player 1:} \\
\texttt{ 1234567} \\
\texttt{╔═══════╗} \\
\texttt{║◌◌◌◌◌◌◌║ 1} \\
\texttt{║◌◌◌◌◌◌◌║ 2} \\
\texttt{║◌TW◌◌┌─╢ 3} \\
\texttt{╟──┐◌◌│H║ 4} \\
\texttt{║C◌├──┼─╢ 5} \\
\texttt{║◌◌│I◌│◌║ 6} \\
\texttt{║◌◌│◌◌│◌║ 7} \\
\texttt{╚══╧══╧═╝} \\
\\ 
\\ 
\texttt{Player 2:} \\
\texttt{ 1234567} \\
\texttt{╔═══════╗} \\
\texttt{║◌◌◌◌TH◌║ 1} \\
\texttt{║◌◌◌◌◌◌◌║ 2} \\
\texttt{║◌◌◌◌◌┌─╢ 3} \\
\texttt{╟──┐◌◌│◌║ 4} \\
\texttt{║I◌├──┼─╢ 5} \\
\texttt{║◌W│◌◌│C║ 6} \\
\texttt{║◌◌│◌◌│◌║ 7} \\
\texttt{╚══╧══╧═╝} \\
\\ 
\texttt{```} \\
            }
        }
    }
    & & \\ \\

    \theutterance \stepcounter{utterance}  
    & & & \multicolumn{2}{p{0.3\linewidth}}{
        \cellcolor[rgb]{0.9,0.9,0.9}{
            \makecell[{{p{\linewidth}}}]{
                \texttt{\tiny{[GM$|$GM]}}
                \texttt{End states:} \\
\texttt{```} \\
\texttt{Player 1:} \\
\texttt{ 1234567} \\
\texttt{╔═══════╗} \\
\texttt{║◌◌◌◌◌◌◌║ 1} \\
\texttt{║◌◌◌H◌◌◌║ 2} \\
\texttt{║◌◌W◌◌┌─╢ 3} \\
\texttt{╟──┐◌T│◌║ 4} \\
\texttt{║C◌├──┼─╢ 5} \\
\texttt{║◌◌│I◌│◌║ 6} \\
\texttt{║◌◌│◌◌│◌║ 7} \\
\texttt{╚══╧══╧═╝} \\
\\ 
\\ 
\texttt{Player 2:} \\
\texttt{ 1234567} \\
\texttt{╔═══════╗} \\
\texttt{║◌◌◌◌◌H◌║ 1} \\
\texttt{║◌◌◌◌◌◌◌║ 2} \\
\texttt{║◌◌◌◌◌┌─╢ 3} \\
\texttt{╟──┐◌◌│◌║ 4} \\
\texttt{║I◌├──┼─╢ 5} \\
\texttt{║TW│◌◌│C║ 6} \\
\texttt{║◌◌│◌◌│◌║ 7} \\
\texttt{╚══╧══╧═╝} \\
\\ 
\texttt{```} \\
            }
        }
    }
    & & \\ \\

    \theutterance \stepcounter{utterance}  
    & & & \multicolumn{2}{p{0.3\linewidth}}{
        \cellcolor[rgb]{0.9,0.9,0.9}{
            \makecell[{{p{\linewidth}}}]{
                \texttt{\tiny{[GM$|$GM]}}
                \texttt{* success: False} \\
\texttt{* lose: False} \\
\texttt{* aborted: True} \\
\texttt{{-}{-}{-}{-}{-}{-}{-}} \\
\texttt{* Shifts: 1.00} \\
\texttt{* Max Shifts: 8.00} \\
\texttt{* Min Shifts: 4.00} \\
\texttt{* End Distance Sum: 19.12} \\
\texttt{* Init Distance Sum: 19.18} \\
\texttt{* Expected Distance Sum: 20.95} \\
\texttt{* Penalties: 13.00} \\
\texttt{* Max Penalties: 12.00} \\
\texttt{* Rounds: 6.00} \\
\texttt{* Max Rounds: 20.00} \\
\texttt{* Object Count: 5.00} \\
            }
        }
    }
    & & \\ \\

\end{supertabular}
}

\end{document}
